% Bagian: sets-functions-relations
% Bab: relations
% Subbagian: special-properties

\documentclass[../../../include/open-logic-section]{subfiles}

\begin{document}

\olfileid[id]{sfr}{rel}{prp}
\olsection{Sifat Khusus Relasi}

\begin{intro}
Beberapa jenis relasi ternyata begitu umum sehingga diberi nama khusus.
Sebagai contoh, $\le$ dan~$\subseteq$ sama-sama menghubungkan domainnya
masing-masing (misalnya, $\Nat$ untuk~$\le$ dan $\Pow{A}$
untuk~$\subseteq$) dengan cara yang serupa. Untuk menentukan dengan tepat
keserupaan dan perbedaan relasi-relasi tersebut, kita mengelompokkannya
menurut beberapa sifat khusus yang dapat dimiliki relasi. Ternyata, beberapa
sifat khusus ini (atau kombinasinya) sangat penting, terutama urutan dan
relasi ekuivalensi.
\end{intro}

\begin{defn}[Refleksivitas]
Suatu relasi $R \subseteq A^2$ bersifat \emph{refleksif} jika dan hanya jika,
untuk setiap $x \in A$, berlaku $Rxx$.
\end{defn}

\begin{defn}[Transitivitas]
Suatu relasi $R \subseteq A^2$ bersifat \emph{transitif} jika dan hanya jika,
setiap kali $Rxy$ dan $Ryz$ berlaku, $Rxz$ juga berlaku.
\end{defn}

\begin{defn}[Simetri]
Suatu relasi~$R \subseteq A^2$ bersifat \emph{simetris} jika dan hanya jika,
setiap kali $Rxy$ berlaku, $Ryx$ juga berlaku.
\end{defn}

\begin{defn}[Antisimetri]
Suatu relasi~$R \subseteq A^2$ bersifat \emph{anti\-si\-metris} jika dan hanya
jika, setiap kali $Rxy$ dan $Ryx$ keduanya berlaku, maka $x=y$ (atau, dengan
kata lain: jika $x\neq y$, maka $\lnot Rxy$ atau $\lnot Ryx$).
\end{defn}

\begin{explain}
Dalam relasi simetris, $Rxy$ dan $Ryx$ selalu berlaku bersama-sama atau
keduanya tidak berlaku. Dalam relasi antisimetris, $Rxy$ dan $Ryx$ hanya dapat
berlaku bersama-sama jika $x = y$. Perhatikan bahwa ini tidak
\emph{mengharuskan} $Rxy$ dan $Ryx$ berlaku ketika $x = y$; hal itu semata-mata
tidak dikesampingkan. Jadi, suatu relasi antisimetris dapat bersifat
refleksif, tetapi tidak setiap relasi antisimetris bersifat refleksif.
Perhatikan pula bahwa bersifat antisimetris berbeda dari sekadar tidak
bersifat simetris. Bahkan, suatu relasi dapat sekaligus bersifat simetris dan
antisimetris (misalnya, relasi identitas).
\end{explain}

\begin{defn}[Keterhubungan]
Suatu relasi $R \subseteq A^2$ bersifat \emph{terhubung} jika untuk semua
$x,y\in A$, apabila $x \neq y$, maka $Rxy$ atau~$Ryx$.
\end{defn}

\begin{prob}
Berikan contoh relasi yang (a) refleksif dan simetris, tetapi tidak transitif;
(b) refleksif dan antisimetris; (c) antisimetris dan transitif, tetapi tidak
refleksif; serta (d) refleksif, simetris, dan transitif. Jangan gunakan relasi
pada bilangan atau himpunan.
\end{prob}

\begin{defn}[Irefleksivitas]
Suatu relasi $R \subseteq A^2$ disebut \emph{irefleksif} jika, untuk semua
$x \in A$, $Rxx$ tidak berlaku.
\end{defn}

\begin{defn}[Asimetri]
Suatu relasi $R \subseteq A^2$ disebut \emph{asimetris} jika tidak ada pasangan
$x,y\in A$ yang memenuhi $Rxy$ dan~$Ryx$ sekaligus.
\end{defn}

Perhatikan bahwa jika $A \neq \emptyset$, maka tidak ada relasi irefleksif
pada~$A$ yang bersifat refleksif, dan setiap relasi asimetris pada~$A$ juga
bersifat antisimetris. Namun, ada $R \subseteq A^2$ yang tidak refleksif dan
juga tidak irefleksif, serta ada relasi antisimetris yang tidak asimetris.

\end{document}
